% 第3章 3.1.2：有限秩截断的算子范数误差与逐点误差（l^2）。
% A(x_n)=(x_n/n)，A_m 保留前 m 个坐标：||A-A_m||=1/(m+1)。
% 坐标投影 P_m：||I-P_m||=1（取 e_{m+1}）；固定 x=(1/n)/||(1/n)||，
% ||(I-P_m)x||=sqrt((pi^2/6-sum_{n<=m}1/n^2)/(pi^2/6))，由 python 直接计算，m=0..20。
\begin{tikzpicture}[font=\figurefont\footnotesize]
\begin{axis}[width=82mm,height=50mm,
  xmin=0,xmax=20.5,ymin=0,ymax=1.12,
  axis x line=bottom,axis y line=left,
  xtick={0,5,10,15,20},ytick={0,0.25,0.5,0.75,1},
  yticklabel style={/pgf/number format/fixed},
  xlabel={保留的坐标数$m$},ylabel={误差},
  xlabel style={font=\figurefont\scriptsize},ylabel style={font=\figurefont\scriptsize},
  tick label style={font=\figurefont\scriptsize,text=FigureInk},
  ymajorgrids,grid style={draw=FigureGrid,line width=.4pt},
  legend style={at={(0.98,0.80)},anchor=east,font=\figurefont\scriptsize,cells={anchor=west}}]
\addplot[FigureYellow,line width=1.1pt,dashed] coordinates {(0,1) (20,1)};
\addplot[FigureGreen,line width=1.1pt,mark=square*,mark size=.9pt] coordinates {(0,1.0000) (1,0.6262) (2,0.4900) (3,0.4154) (4,0.3668) (5,0.3320) (6,0.3055) (7,0.2845) (8,0.2673) (9,0.2529) (10,0.2405) (11,0.2298) (12,0.2205) (13,0.2122) (14,0.2047) (15,0.1980) (16,0.1919) (17,0.1864) (18,0.1813) (19,0.1765) (20,0.1722)};
\addplot[FigureBlue,line width=1.1pt,mark=*,mark size=1pt] coordinates {(0,1.0000) (1,0.5000) (2,0.3333) (3,0.2500) (4,0.2000) (5,0.1667) (6,0.1429) (7,0.1250) (8,0.1111) (9,0.1000) (10,0.0909) (11,0.0833) (12,0.0769) (13,0.0714) (14,0.0667) (15,0.0625) (16,0.0588) (17,0.0556) (18,0.0526) (19,0.0500) (20,0.0476)};
\legend{$\|I-P_m\|=1$,$\|(I-P_m)x\|$（固定$x$）,$\|A-A_m\|=1/(m+1)$}
\end{axis}
\end{tikzpicture}
